2. \underline{意大利学派的相交理论}

设 $S'$ 与 $S$ 是域 $k$ 上的两个光滑射影曲面，$f:S'\longrightarrow S$ 是双有理态射。若 $D$ 与 $E$ 是 $S$ 上的两个除子，则

\begin{equation*}
(D,E)=(f^*D,f^*E),
\end{equation*}

%%%%%%%%%% source scan idx 19 | folio 12 %%%%%%%%%%

例如，移动 $D$ 与 $E$ 即可看出这一点。

因此，自然要寻找相交数的一个“双有理”表达式。我将说明意大利学派的做法。主要参考文献是 [1]，其中附有很长的参考文献表。另见 Manin [6] 及 Northcott [12]、[13]。

(2.0) 设 $S$ 是二维正则局部环的谱，$s$ 是 $S$ 的闭点，$k$ 是 $k(s)$ 的子域，且 $k(s)$ 在其上有限。相交数和 Euler--Poincaré 特征数均相对于 $k$ 计算。

本节中，若 $S$-概形 $p:S'\longrightarrow S$ 正则且连通，$p$ 固有且诱导同构 $S'\setminus p^{-1}(s)\longrightarrow S\setminus\{s\}$，就称它为 $S$ 的\underline{双有理变换}。关于 $S$ 的所有论述，也适用于其双有理变换在闭点处的局部化。

(2.1) 设 $f:(S'',p'')\longrightarrow(S',p')$ 是 $S$ 的双有理变换之间的 $S$-概形态射。已知 [5]，$f$ 有唯一的分解

\begin{equation*}
f=f_n\cdots f_1,\qquad f_i:S_{i+1}\longrightarrow S_i,
\end{equation*}
使得 $S_{i+1}$ 由 $S_i$ 沿闭点的有限集合 $X_i$ 爆破而得，且 $f_i(X_{i+1})\subset X_i$。集合 $X_i$ 恰由 $S_i$ 中使得 $S''$ 到 $S_i$ 的投影不为同构的点组成。

特别地，$f$ 在 $S'$ 的某个有限子集的补集上诱导同构，该补集包含 $p'^{-1}(s)$ 各不可约分支的泛点。因此 $f^{-1}$ 定义单射

%%%%%%%%%% source scan idx 20 | folio 13 %%%%%%%%%%

将 $p'^{-1}(s)$ 的不可约分支的集合 $\operatorname{Irr}(p'^{-1}(s))$ 映入 $\operatorname{Irr}(p''^{-1}(s))$。

$S$ 的双有理变换没有非平凡的 $S$-自同构。用 $I$ 表示 $S$ 的双有理变换的 $S$-同构类集合，按控制关系排序。集合 $I$ 向右滤过，且有最小元 $S$。

\underline{定义} 2.2. \underline{$S$ 中与 $s$ 无限邻近的点的集合 $S^*$ 定义为正向极限集合}

\begin{equation*}
S^*=\varinjlim_I\operatorname{Irr}(p^{-1}(s)) \quad .
\end{equation*}
若元素 $t\in S^*$ 在 $S$ 的双有理变换 $S'$ 上由完备约化不可约曲线 $T$ 表示，则域 $H^0(T,\underline{O}_T)$ 只依赖于 $t$，与 $S'$ 的选择无关。称其为 $t$ 的\underline{剩余域} $k(t)$。

设 $S_1$ 是 $S$ 沿 $s$ 的爆破。$\{s\}$ 在 $S_1$ 中的逆像射影直线 $s^*$ 称为 $s$ 的\underline{小曲线（curvetta）}，它定义 $S^*$ 的元素 $s^*$，其剩余域为 $k(s)$。

设 $S'$ 是 $S$ 的双有理变换，$t$ 是 $S'$ 的闭点，$S'_t$ 是 $S'$ 在 $t$ 处的局部环的谱。有一个显然的单射 $(S'_t)^*\hookrightarrow S^*$；特别地，$t$ 定义 $S^*$ 中具有相同剩余域的元素 $t^*$。用 $(S')^*$ 表示当 $t$ 遍历 $S'$ 的闭点时，各 $(S'_t)^*$ 在 $S^*$ 中的并。

由 2.1 中对 $S$ 的双有理变换的描述，得到下述命题。

%%%%%%%%%% source scan idx 21 | folio 14 %%%%%%%%%%

\underline{命题} 2.3. \underline{对每个 $t\in S^*$，存在唯一的一列 $S$ 的双有理变换 $(S_i)_{0\leq i\leq n+1}$ 及闭点 $t_i\in S_i$ $(0\leq i\leq n)$，使得}

(i)\quad $S_0=S$，\underline{且 $S_{i+1}$ 是 $S_i$ 沿 $\{t_i\}$ 的爆破；}

(ii)\quad $t_{i+1}\in t_i^*$ $(0\leq i<n)$，\underline{且 $t=t_n^*$。}

\underline{此外，对 $S$ 的任一双有理变换 $(S',p')$，元素 $t$ 属于 $\operatorname{Irr}(p'^{-1}(s))$ 在 $S^*$ 中的像，当且仅当 $S'$ 控制 $S_{n+1}$。}

整数 $n$ 称为 $t$ 的阶。用 $S(t)$（分别用 $S_+(t)$）表示 $S$ 的双有理变换 $S_n$（分别为 $S_{n+1}$）。沿用记号，也用 $t$ 表示 $S_n$ 的闭点 $t_n$，用 $t^*$ 表示它在 $S_{n+1}$ 中的逆像 $t_n^*$。

赋予 $S^*$ 如下序关系：

\begin{equation*}
t_1\prec t_2
\quad\Longleftrightarrow\quad
S_+(t_2)\ \text{控制}\ S_+(t_1) \quad .
\end{equation*}
若 $t_1\prec t_2$，则 $k(t_2)$ 是 $k(t_1)$ 的扩张。有序集 $S^*$ 有最小元 $s^*$，且 $S^*$ 的每个有界子集均有限且全序（$S^*$ 是“以 $s^*$ 为根的树”）。此外，$t\in S^*$ 的阶 $n$ 是树 $S^*$ 中从 $s^*$ 到 $t$ 的距离。

下述命题使 2.1 中对 $S$ 的双有理变换的描述具体化。

\underline{命题} 2.4. \underline{将双有理变换 $p:S'\to S$ 对应到 $S^*$ 的子集 $\operatorname{Irr}(p^{-1}(s))$ 的映射，诱导 $I$（2.1）与 $S^*$ 的有限子集之间的双射，这些子集具有如下性质：随每个元素一起，}

%%%%%%%%%% source scan idx 22 | folio 15 %%%%%%%%%%

\underline{包含它的所有前驱。}

设 $D$ 是 $S$ 上的除子 $\geq0$，$t\in S^*$。若 $S(t)$ 的闭点 $t$（2.3）属于 $D$ 在 $S(t)$ 上的严格变换（即 $D-\{s\}$ 的概形闭包），就说 $t$ \underline{属于} $D$。

\underline{定义} 2.5. \underline{设 $t,t'\in S^*$。若 $t'\in S_+(t)^*$（2.2），且属于 $t$ 的小曲线 $t^*$（2.3），就说 $t'$ 邻接于 $t$（proximate）。}

若 $t'$ 是 $t$ 的直接后继，则 $t'$ 邻接于 $t$。若在序列 $(t_i)_{0\leq i\leq n}$ 中每个 $t_{i+1}$ 都是 $t_i$ 的直接后继，且 $t_n$ 邻接于 $t_0$，则每个满足 $0<i\leq n$ 的 $t_i$ 都邻接于 $t_0$，而在 $0<i\leq n-1$ 时 $t_n$ 不邻接于 $t_i$。

(2.6) 设 $A$ 是维数为 $n$ 的 Noether 局部环，$T$ 是其谱，$m$ 是其极大理想，$k=A/m$ 是其剩余域，$P=\sum_{0\leq i<n}a_iX^i$ 是其 Hilbert--Samuel 多项式：

\begin{equation*}
P(r)=\dim_k(m^r/m^{r+1})\qquad\text{对充分大的 }r.
\end{equation*}
回顾闭点 $s$ 在 $T$ 中的重数 $\mu_s$ 的定义：

\begin{equation*}
a_{n-1}=\frac{\mu_s}{(n-1)!} \quad .
\end{equation*}
当 $n=1$ 时，这就是说

\begin{equation*}
\mu_s=\dim_k(m^r/m^{r+1})\qquad\text{对充分大的 }r.
\end{equation*}
从几何上说，$\mu_s$ 可定义为 $T$ 沿闭子概形 $\{s\}$ 的爆破的特殊纤维在射影空间 $\mathbf P_k(m/m^2)$ 中的次数。

%%%%%%%%%% source scan idx 23 | folio 16 %%%%%%%%%%

回顾：若 $Y$ 与 $Z$ 是概形 $X$ 的两个闭子概形，则 $Y$ 在 $X$ 沿 $Z$ 的爆破 $\widetilde X$ 中的严格变换——即 $Y\setminus Z$ 在 $\widetilde X$ 中的概形闭包——恰为 $Y$ 沿 $Y\cap Z$ 的爆破。

假设 $T$ 是正则概形 $U$ 的闭子概形。设 $p:U_1\longrightarrow U$ 是 $U$ 沿 $\{s\}$ 的爆破，$s^*$ 是 $s$ 在 $U_1$ 中的逆像，$T'$ 是 $T$ 在 $U_1$ 中的严格变换。上述讨论表明，$T$ 在 $s$ 处的重数是子概形 $s^*\cap T'$ 在射影空间 $s^*$ 中的次数。

再假定 $T$ 是 $U$ 上的除子。设 $T^*$ 是其在 $U_1$ 中的全变换（逆像除子），由 $T^*=\nu s^*+T'$ 定义 $\nu$。若 $X$ 是 $s^*$ 中的直线，且 Euler--Poincaré 特征数作为 $k$ 上维数的交错和计算，则

\begin{equation}\tag{2.6.1}
(X,s^*)=
\chi(\underline{O}_X\overset{\mathbb L}{\otimes}\underline{O}_{s^*})
=\deg_X\underline{O}(s^*)=-1,
\end{equation}
并且我们已经看到

\begin{equation}\tag{2.6.2}
(X,T')=
\chi(\underline{O}_X\overset{\mathbb L}{\otimes}\underline{O}_{T'})
=\deg_{s^*}(s^*\cap T')=\mu_s \quad .
\end{equation}
投影公式

\begin{equation*}
Rp_*(\underline{O}_X\overset{\mathbb L}{\otimes}\underline{O}_{T^*})
=Rp_*(\underline{O}_X\overset{\mathbb L}{\otimes}Lp^*\underline{O}_T)
=Rp_*\underline{O}_X\overset{\mathbb L}{\otimes}\underline{O}_T
\end{equation*}
及其推论

\begin{equation*}
(X,T^*)=\nu(X,s^*)+(X,T')=0
\end{equation*}
结合（2.6.1）和（2.6.2）给出恒等式

\begin{equation}\tag{2.6.3}
\mu_s=\nu \quad .
\end{equation}

%%%%%%%%%% source scan idx 24 | folio 17 %%%%%%%%%%

\underline{命题} 2.7. \underline{设 $T$ 是正则概形 $U$ 上的除子 $\geq0$，$s$ 是 $U$ 的闭点，$U_1$ 是 $U$ 沿 $\{s\}$ 的爆破，$s^*$ 是 $s$ 在 $U_1$ 中的逆像，$T'$ 是 $T$ 在 $U_1$ 上的严格变换，$T^*$ 是它的全变换。}

\underline{下列三个量相等：}

(i)\quad \underline{点 $s$ 在 $T$ 上的 Samuel 意义下的重数；}

(ii)\quad \underline{$s^*\cap T'$ 在射影空间 $s^*$ 中的次数；}

(iii)\quad \underline{$s^*$ 在除子 $T^*$ 中的重数。}

2.8. 回到 2.0 的假设。设 $D$ 是 $S$ 上的除子 $\geq0$，$t\in S^*$。$D$ 在 $t$ 处的\underline{重数} $\mu_t(D)$，定义为 $D$ 在 $S(t)$（2.3）上的严格变换在 $t$ 处的重数。函数 $t\mapsto\mu_t(D)$ 的支集是属于 $D$ 的与 $s$ 无限邻近的点的集合（2.4）。当 $D$ 正则时，各数 $\mu_t(D)[k(t):k(s)]$ 为零或一：事实上，$D$ 的严格变换 $D'$ 同构于 $D$，再应用重数的定义（2.6）即可。

相交数可以用重数来计算。

\underline{定理} 2.9. \underline{若除子 $D$ 与 $E$ 只在 $s$ 相交，则}

\begin{equation*}
(D,E)=\sum_{t\in S^*}\mu_t(D)\mu_t(E)[k(t):k],
\end{equation*}
\underline{其中对既属于 $D$ 又属于 $E$ 的与 $s$ 无限邻近的点求和。}